% Generated by report/tools/tables.py from the result files of the repository. Do not edit.
\begin{table}[tbp]
\centering
\footnotesize
\setlength{\tabcolsep}{3pt}
\caption[Pineau's grid: 2018 against classical RL 2026]{Pineau's grid (criteria Q2 to Q8) and the X criteria drawn from her ICLR 2018 talk, read with the same grid and the same readers on two samples: our reading of 50 papers from 2018 (ICML, NeurIPS, ICLR) and 50 classical RL papers from 2026 (14 from the first draw, 36 from the A3 queue). This is the main comparison, the only one a verdict rests on (deviation 20 of \file{study/PROTOCOL.md}). $k/n$: papers that meet the criterion out of the papers where it applies (papers with experiments, Q1 not `no'); the percentage carries its 95\% Wilson interval. Gap: classical RL 2026 minus 2018, in percentage points, with the 95\% Newcombe interval; $p$: two-sided Fisher exact test; Holm $p$: Holm correction over the eleven criteria of grid v1. \textbf{In bold}, the criteria whose gap passes the Holm correction (Holm $p < 0.05$). For X3, a lower share is progress; Q2 (neural networks) describes the population studied, not a reproducibility practice. Pineau's own figures are in Table~\ref{tab:sondage-pineau-contexte}. Source: \file{study/extension\_stats.json} (script \file{study/extension\_stats.py}, written before any reading of the added papers).}
\label{tab:sondage-pineau}
\begin{tabularx}{\linewidth}{@{}l>{\raggedright\arraybackslash}Xr r@{\hspace{0.28em}}l r r@{\hspace{0.28em}}l r@{\hspace{0.28em}}l rr@{}}
\toprule
 & & \multicolumn{3}{c}{\makecell{Our reading\\{}of 2018}} & \multicolumn{3}{c}{\makecell{Classical RL\\{}2026}} & \multicolumn{2}{c}{Gap} & &  \\
\cmidrule(lr){3-5}\cmidrule(lr){6-8}
Code & Criterion & $k/n$ & \multicolumn{2}{c}{\% [Wilson]} & $k/n$ & \multicolumn{2}{c}{\% [Wilson]} & \multicolumn{2}{c}{[Newcombe]} & $p$ & Holm $p$ \\
\midrule
\textbf{Q2} & \textbf{Neural networks used} & \textbf{38/48} & \textbf{79} & \textbf{[66,\,88]} & \textbf{49/50} & \textbf{98} & \textbf{[90,\,100]} & \textbf{+19} & \textbf{[+6,\,+32]} & \textbf{0.003} & \textbf{0.03} \\
Q3 & All of the method's hyperparameters & 26/47 & 55 & [41,\,69] & 36/50 & 72 & [58,\,83] & +17 & [−2,\,+34] & 0.10 & 0.58 \\
Q4 & All of the baselines' hyperparameters & 17/48 & 35 & [23,\,50] & 14/49 & 29 & [18,\,42] & −7 & [−25,\,+11] & 0.52 & 1 \\
\textbf{Q5} & \textbf{Code linked} & \textbf{15/48} & \textbf{31} & \textbf{[20,\,45]} & \textbf{37/50} & \textbf{74} & \textbf{[60,\,84]} & \textbf{+43} & \textbf{[+23,\,+58]} & \textbf{4\,×\,10\textsuperscript{−5}} & \textbf{0.0004} \\
Q6 & Hyperparameter selection described & 12/47 & 26 & [15,\,40] & 24/50 & 48 & [35,\,61] & +22 & [+3,\,+39] & 0.03 & 0.24 \\
Q7 & Evaluation on held-out conditions & 16/48 & 33 & [22,\,47] & 22/49 & 45 & [32,\,59] & +12 & [−8,\,+30] & 0.30 & 1 \\
Q8 & Statistical test applied & 3/48 & 6 & [2,\,17] & 6/50 & 12 & [6,\,24] & +6 & [−7,\,+18] & 0.49 & 1 \\
\textbf{X2} & \textbf{Uncertainty reported} & \textbf{30/48} & \textbf{63} & \textbf{[48,\,75]} & \textbf{44/50} & \textbf{88} & \textbf{[76,\,94]} & \textbf{+26} & \textbf{[+8,\,+41]} & \textbf{0.005} & \textbf{0.04} \\
X3 & Only the best runs kept & 9/48 & 19 & [10,\,32] & 7/50 & 14 & [7,\,26] & −5 & [−20,\,+10] & 0.59 & 1 \\
X4 & Baselines tuned with the same care & 6/48 & 13 & [6,\,25] & 8/49 & 16 & [9,\,29] & +4 & [−11,\,+18] & 0.77 & 1 \\
\textbf{X7} & \textbf{Compute quantified} & \textbf{13/48} & \textbf{27} & \textbf{[17,\,41]} & \textbf{29/50} & \textbf{58} & \textbf{[44,\,71]} & \textbf{+31} & \textbf{[+11,\,+47]} & \textbf{0.002} & \textbf{0.02} \\
\bottomrule
\end{tabularx}
\end{table}
